% Complete Section 8.8, original PDF 117-122 / printed 107-112.
% block: elasticity-heading page:117 type:section
\section{{弹性力学}}\label{sec:elasticity}

% block: elasticity-feec-introduction page:117 type:ordinary-paragraph
{FEEC 观点的显著成功之一, 是发展了弹性方程的稳定混合有限元方法. 此时相关的 Hilbert 复形不是 de Rham 复形, 而是某种版本的\emph{弹性复形}.}

% block: elasticity-equilibrium-constitutive-laws page:117 type:ordinary-paragraph
{设 $\Omega$ 是 $\mathbb R^n$ 中由弹性体占据的区域. 弹性力学中的主要未知量是位移场 $u:\Omega\to\mathbb R^n$ 和应力场 $\sigma:\Omega\to\mathbb S:=\mathbb R_{\mathrm{symm}}^{n\times n}$. 在稳态下, 应力场必须满足平衡条件}
\[
\operatorname{div}\sigma=f,
\]
{其中 $f:\Omega\to\mathbb R^n$ 是载荷向量场, 给出所施加的单位体积体力. 此外, 位移与应力通过本构方程联系}
\[
A\sigma=\epsilon u,
\]
{其中 $\epsilon$ 表示无穷小应变算子, 即梯度的对称部分, 柔度张量 $A$ 是可能变化的对称正定映射 $\mathbb S\to\mathbb S$. 加上适当边界条件, 例如 $\partial\Omega$ 上的位移条件 $u=0$, 就得到完整的弹性边值问题.}

% block: elasticity-divergence-hilbert-complex pages:117,118 type:ordinary-paragraph
{逐行作用的散度算子 $\operatorname{div}$, 定义了从 $L^2$ 对称张量场 (即对称 $n\times n$ 矩阵值函数) 到 $L^2$ 向量场的闭算子. 其定义域是由散度属于 $L^2$ 的 $L^2$ 对称张量场组成的空间 $H(\operatorname{div},\Omega;\mathbb S)$. 在 $L^2$ 对称张量场空间上, 用柔度张量 $A$ 给 $L^2$ 内积加权, 因而得到 Hilbert 复形的一个片段}
\begin{equation}\label{eq:elasticity-divergence-complex}
L_A^2(\Omega;\mathbb S)\xrightarrow{(\operatorname{div},H(\operatorname{div};\mathbb S))}L^2(\Omega;\mathbb R^n)\to0.
\end{equation}
{它完全类似于复形 \eqref{eq:weighted-divergence-complex}, 区别在于这里的散度算子将对称张量映为向量, 而非将向量映为标量. 对称性约束使弹性系统的离散化困难得多.}

% block: elasticity-mixed-weak-formulation page:118 type:ordinary-paragraph
{与处理 \eqref{eq:weighted-divergence-complex} 时一样, 立即得到此复形的抽象 Hodge Laplace 问题的弱形式. 求 $\sigma\in H(\operatorname{div};\mathbb S)$, $u\in L^2(\Omega;\mathbb R^n)$, 使得}
\[
\begin{aligned}\langle A\sigma,\tau\rangle-\langle u,\operatorname{div}\tau\rangle&=0,&\tau&\in H(\operatorname{div},\Omega;\mathbb S),\\\langle\operatorname{div}\sigma,v\rangle&=\langle f,v\rangle,&v&\in L^2(\Omega;\mathbb R^n).\end{aligned}
\]
{显然, 第一个方程是本构方程及位移边界条件的弱形式, 第二个方程则直接表达平衡条件. 下一个定理断言, 作用于对称张量的散度算子实际上是满射. 这意味着 Hilbert 复形 \eqref{eq:elasticity-divergence-complex} 是闭的, 且没有调和形式. 于是由 Chapter~\ref{chap:hilbert-complexes} 的一般理论可得, 混合弹性系统适定.}

% block: symmetric-tensor-divergence-surjectivity page:118 type:theorem

\begin{Theorem}\label{thm:symmetric-tensor-divergence-surjectivity}
{散度算子 (逐行作用) 将 $H(\operatorname{div},\Omega;\mathbb S)$ 满射到 $L^2(\Omega;\mathbb R^n)$.}
\end{Theorem}

% block: skew-tensor-space page:118 type:ordinary-paragraph
{事实上, 将证明一个更强的结果. 类似于对称 $n\times n$ 张量场空间 $\mathbb S$, 以 $\mathbb K$ 表示反对称 $n\times n$ 张量场空间. 它们的维数分别为 $\dim\mathbb S=\binom{n+1}2$, $\dim\mathbb K=\binom n2$. 如常, 对任意矩阵 $\tau\in\mathbb R^{n\times n}$, $\operatorname{skw}\tau\in\mathbb K$ 表示其反对称部分.}

% block: divergence-skew-surjectivity page:118 type:theorem

\begin{Theorem}\label{thm:divergence-skew-surjectivity}
{给定 $f\in L^2(\Omega;\mathbb R^n)$ 和 $g\in L^2(\Omega;\mathbb K)$, 存在 $\tau\in H(\operatorname{div},\Omega;\mathbb R^{n\times n})$, 使得}
\[
\operatorname{div}\tau=f,\qquad\operatorname{skw}\tau=g.
\]
\end{Theorem}

% block: symmetric-divergence-corollary page:118 type:ordinary-paragraph
{在 Theorem~\ref{thm:divergence-skew-surjectivity} 中取 $g=0$, 即得 Theorem~\ref{thm:symmetric-tensor-divergence-surjectivity}.}

% block: divergence-skew-proof-first-ingredient page:118 type:proof
\begin{Proof}
{只在 $n=2$ 或 $3$ 维中证明此结果, 在证明末尾对一般情形作些说明.}

{构造基于两个要素. 首先是作用于 $H^1(\Omega;\mathbb R^n)$ 的散度的满射性: 给定 $q\in L^2(\Omega)$, 可以找到 $v\in H^1(\Omega;\mathbb R^n)$, 使 $\operatorname{div}v=q$. 前一节关于 Stokes 方程的讨论中, 研究了作用于 $\mathring H^1(\Omega;\mathbb R^n)$ 的散度满射性这一更困难的结果. 没有边界条件时, 结果容易得多. 例如, 可以将 $q$ 零延拓至一个包含 $\Omega$ 的球, 然后在球上求解强迫项为 $q$ 的 Dirichlet 问题, 得到一个 $H^2$ 函数, 最后取 $v$ 为该解的梯度在 $\Omega$ 上的限制.}

% block: divergence-skew-proof-two-dimensional-identity pages:118,119 type:ordinary-paragraph
{第二个要素是一个初等微积分恒等式. 主要出于记号上的原因, 先在 $n=2$ 维中单独介绍, 再介绍三维情形. 当 $n=2$ 时, 空间 $\mathbb K$ 是一维的. 因而可用属于 $L^2(\Omega)$ 的标量函数 $g$ 来陈述定理, 且必须证明存在 $\tau$, 使 $\operatorname{div}\tau=f$ 且 $\operatorname{asym}\tau=g$, 其中 $\operatorname{asym}\tau:=\tau_{12}-\tau_{21}$. 在二维情形中, 所需的初等恒等式是: 二维向量场的矩阵旋度的不对称部分, 等于其散度的负值. 即}
\begin{equation}\label{eq:asym-curl-divergence-identity}
\operatorname{asym}\operatorname{curl}\psi=-\operatorname{div}\psi,\qquad\psi\in H^1(\Omega;\mathbb R^2).
\end{equation}
{(这里旋度算子是作用于 $\psi$ 各标量分量的二维向量旋度, 得到 $2\times2$ 矩阵的各行.) 此恒等式也由下图的交换性表达:}
\[
\begin{tikzcd}[column sep=large,row sep=large]
H^1(\Omega;\mathbb R^2)\arrow[r,"\operatorname{curl}"]\arrow[d,"id"']&L^2(\Omega;\mathbb R^{2\times2})\arrow[d,"-\operatorname{asym}"]\\
H^1(\Omega;\mathbb R^2)\arrow[r,"\operatorname{div}"']&L^2(\Omega)
\end{tikzcd}
\]

% block: divergence-skew-two-dimensional-construction page:119 type:ordinary-paragraph
{有了这两个要素, 就可进行构造. 对任意 $f\in L^2(\Omega;\mathbb R^2)$, 可找到 $\rho\in H^1(\Omega;\mathbb R^{2\times2})$, 使得}
\begin{equation}\label{eq:elasticity-divergence-lifting}
\operatorname{div}\rho=f.
\end{equation}
{事实上, 由于没有对 $f$ 施加对称性约束, 可以直接利用散度在 $H^1(\Omega;\mathbb R^2)$ 上的满射性, 独立地定义 $\rho$ 的两行. 接着构造标量函数 $-g+\operatorname{asym}\rho$, 再次应用散度的满射性, 得到向量场 $\psi\in H^1(\Omega;\mathbb R^2)$ 满足}
\begin{equation}\label{eq:elasticity-skew-correction-lifting}
\operatorname{div}\psi=-g+\operatorname{asym}\rho.
\end{equation}
{注意, $\psi$ 的矩阵旋度属于 $L^2$ 且无散, 因而属于 $H(\operatorname{div},\Omega;\mathbb R^{2\times2})$. 因此}
\[
\tau:=\rho+\operatorname{curl}\psi\in H(\operatorname{div},\Omega;\mathbb R^{2\times2}).
\]
{最后, 由 \eqref{eq:elasticity-divergence-lifting}, \eqref{eq:asym-curl-divergence-identity} 和 \eqref{eq:elasticity-skew-correction-lifting}, 得到}
\[
\begin{aligned}
\operatorname{div}\tau&=\operatorname{div}\rho=f,\\
\operatorname{asym}\tau&=\operatorname{asym}\rho+\operatorname{asym}\operatorname{curl}\psi=\operatorname{asym}\rho-\operatorname{div}\psi=g,
\end{aligned}
\]
{如所要求. 这就证明了二维情形的 Theorem~\ref{thm:divergence-skew-surjectivity}.}

% block: divergence-skew-three-dimensional-identity page:119 type:ordinary-paragraph
{在三维中, 空间 $\mathbb K$ 本身为三维, 即 $\mathbb K\cong\mathbb R^3$. 因而取 $g$ 为三维向量场, 并施加约束 $\operatorname{asym}\tau=g$, 其中现在}
\[
\operatorname{asym}\tau=(\tau_{23}-\tau_{32},\tau_{31}-\tau_{13},\tau_{12}-\tau_{12}).
\]
{还需要算子 $S:\mathbb R^{3\times3}\to\mathbb R^{3\times3}$, 定义为}
\[
S\rho=\rho^T-(\operatorname{Tr}\rho)I,
\]
{其中 $\operatorname{Tr}$ 表示矩阵的迹, $I$ 表示 $3\times3$ 单位矩阵. 此算子可逆, 且}
\[
S^{-1}\mu=\mu^T-\frac12(\operatorname{Tr}\mu)I.
\]
{现在, 三维中的相应恒等式可表达为下图的交换性}
\[
\begin{tikzcd}[column sep=large,row sep=large]
H^1(\Omega;\mathbb R^{3\times3})\arrow[r,"\operatorname{curl}"]\arrow[d,"S"']&L^2(\Omega;\mathbb R^{3\times3})\arrow[d,"-\operatorname{asym}"]\\
H^1(\Omega;\mathbb R^{3\times3})\arrow[r,"\operatorname{div}"']&L^2(\Omega;\mathbb R^3)
\end{tikzcd}
\]

% block: divergence-skew-three-dimensional-construction page:120 type:ordinary-paragraph
{接着论证与二维情形一样进行. 先定义 $\rho\in H^1(\Omega;\mathbb R^{3\times3})$, \newline 再定义 $\psi\in H^1(\Omega;\mathbb R^{3\times3})$, 使得}
\[
\operatorname{div}\rho=f,\qquad\operatorname{div}\psi=-g+\operatorname{asym}\rho,
\]
{然后令}
\[
\tau=\rho+\operatorname{curl}S^{-1}\psi.
\]
\end{Proof}

% block: elasticity-surjectivity-general-dimension page:120 type:remark

\begin{Remark}\label{rem:elasticity-general-dimension}
{为在一般维数 $n$ 中证明上述定理, 最好将定理陈述及证明中涉及的空间, 即标量值, 向量值和矩阵值函数空间, 视为微分形式空间. 这样得到的微分形式不是标量值的, 而是取值于 $\operatorname{Alt}^m\mathbb R^n$, 其中 $m=n-1$ 或 $n-2$. 因而, 例如给定 $f$ 所在的空间 $L^2(\Omega;\mathbb R^n)$, 实际上应视为 $L^2\Lambda^n(\Omega)\otimes\operatorname{Alt}^{n-1}\mathbb R^n$, 即 $\Omega$ 上取值于 $\operatorname{Alt}^{n-1}\mathbb R^n$ 的 $L^2$ 微分 $n$-形式空间. 回顾 $\operatorname{Alt}^{n-1}\mathbb R^n$ 与 $\mathbb R^n$ 典范同构 (向量代理), 且微分 $n$-形式可视为普通标量值函数 (标量代理), 因而确实可将 $L^2\Lambda^n(\Omega)\otimes\operatorname{Alt}^{n-1}\mathbb R^n$ 的元素与 $\Omega$ 上的 $\mathbb R^n$ 值 $L^2$ 函数典范等同. 类似地, 可等同给定 $g$ 所在的空间以及定义 $\rho$, $\psi$ 和 $\tau$ 所在的空间. 在此背景下, 可以对一般的 $n$ 定义算子 $\operatorname{asym}$ 和 $S$ 的对应物, 并建立证明所需的交换图.}
\end{Remark}

% block: two-dimensional-elasticity-complex page:120 type:ordinary-paragraph
{现在转向弹性系统的离散化, 如前所见, 它就是复形 \eqref{eq:elasticity-divergence-complex} 的抽象 Hodge Laplace 问题. 因而需要发展 \eqref{eq:elasticity-divergence-complex} 中各空间的有限元子复形, 分别用于离散化应力和位移. 如上所述, 对称性约束使此事相当困难. 将序列 \eqref{eq:elasticity-divergence-complex} 延长为更长的复形是有用的, 就像 \eqref{eq:weighted-divergence-complex} 延长为 de Rham 复形一样. 在二维中, 这个更长的\emph{弹性复形}可写成}
\begin{equation}\label{eq:two-dimensional-elasticity-complex}
0\to L^2(\Omega)\xrightarrow{(\operatorname{curl}\operatorname{curl},H^2)}L^2(\Omega;\mathbb S)\xrightarrow{(\operatorname{div},H(\operatorname{div};\mathbb S))}L^2(\Omega;\mathbb R^2)\to0,
\end{equation}
{其中为简单起见, 略去加权系数 $A$. 这里算子 $\operatorname{curl}\operatorname{curl}$ 是先对标量函数取向量值旋度, 再对向量函数取矩阵值旋度所得到的复合. 所得张量场,}
\[
\operatorname{curl}\operatorname{curl}\phi=
\begin{pmatrix}\partial^2\phi/\partial x_2^2&-\partial^2\phi/\partial x_2^2\\-\partial^2\phi/\partial x_2^2&\partial^2\phi/\partial x_1^2\end{pmatrix},
\]
{在弹性理论中称为与标量势函数 $\phi$ 关联的 Airy 应力函数.}

% block: three-dimensional-elasticity-complex pages:120,121 type:ordinary-paragraph
{在三维中, 弹性复形为}
\[
\begin{aligned}
0\to L^2(\Omega;\mathbb R^3)&\xrightarrow{(\epsilon,H^1(\Omega;\mathbb R^3))}L^2(\Omega;\mathbb S)\\
&\xrightarrow{\operatorname{curl}T\operatorname{curl}}L^2(\Omega;\mathbb S)\xrightarrow{(\operatorname{div},H(\operatorname{div};\mathbb S))}L^2(\Omega;\mathbb R^3)\to0.
\end{aligned}
\]
{此时算子 $\operatorname{curl}T\operatorname{curl}$ 称为 St. Venant 算子或不相容算子, 它通过对对称 $3\times3$ 矩阵场逐行取旋度, 再转置, 再取旋度而得到 (或等价地, 先逐行再逐列取旋度). 其定义域由 $\operatorname{curl}T\operatorname{curl}$ 仍属于 $L^2$ 的 $L^2$ 对称矩阵场组成. 注意, 它的像中的矩阵场是对称的. \BookCite{2} 对这一不相容算子作了广泛分析.}

% block: symmetric-elasticity-finite-elements page:121 type:ordinary-paragraph
{若交换 Airy 应力函数 \eqref{eq:maxwell-faraday} 的对角元素, 并将非对角元素变号, 就得到 $\phi$ 的 Hessian. 这表明二维弹性复形的离散化与 $H^2$ 标量有限元密切相关. 这一观察帮助引导了二维中首批使用多项式试函数的稳定弹性混合有限元的发现. 它们于 2002 年在 \BookCite{15} 中提出. 应力单元和位移单元与 $H^2$ 的经典 Hermite 五次, 或 Argyris, 有限元子空间一起, 构成 \eqref{eq:two-dimensional-elasticity-complex} 的有限元子复形. 此后提出了许多替代方案和改进. \BookCite{15} 的单元先在最低次情形 \BookCite{1}, 后在一般次数情形 \BookCite{8}, 被推广至三维. \BookCite{53} 和 \BookCite{54} 发展了适用于任意维数的更高效单元. 后一文献提供了相关文献的更广泛综述.}

% block: weak-symmetry-elasticity-complex page:121 type:ordinary-paragraph
{简要提及发展高效弹性混合有限元的另一种方法, 即弱施加对称性 (这一思想可追溯到有限元早期 \BookCite{37}). 该方法得到不同的混合弱形式, 在不施加对称性的空间中求应力场, 改以应力的反对称部分与 Lagrange 乘子正交来弱施加对称性. 所得系统仍可视为某个复形的抽象 Hodge Laplacian, 但相关复形此时不再是 \eqref{eq:weighted-divergence-complex}, 而是}
\[
L_A^2(\Omega;\mathbb R^{n\times n})\xrightarrow{((\operatorname{div},\operatorname{skw}),H(\operatorname{div};\mathbb R^{2\times2}))}L^2(\Omega;\mathbb R^n)\times L^2(\Omega;\mathbb K)\to0.
\]
{注意, 此时的微分由一对算子组成: 它将矩阵场同时映为其散度和反对称部分. 上面已在 Theorem~\ref{thm:divergence-skew-surjectivity} 中证明此微分是满射, 因而再次得到适定的混合形式. 此片段也可延长成更长的复形 (\emph{弱对称弹性复形}). 离散化此复形可采用比标准弹性复形所需单元更简单的单元. \BookCite{12} 在三维中发展了这样的单元, 此外也提出了各种替代单元. 例如, 见 \BookCite{32}, \BookCite{46} 以及 \BookCite{47} 引用的工作.}
% Original PDF 122 is a blank production leaf, retained unchanged in the source PDF.
